\subsubsection{სისტემის ოპტიმიზაცია}

პროექტის მეორე ნახევარში ორი დიდი ოპტიმიზაცია განვახორციელეთ: პოზიციური
B+ ინდექსი და ნაგვის შემგროვებელი სისტემა. ორივე იმ პრობლემებს პასუხობს,
რომლებიც ყველა სერიოზულ CRDT იმპლემენტაციას აწუხებს და რომელთა
გადაწყვეტასაც Yjs \cite{yjs} საკუთარ დიზაინში დიდ ადგილს უთმობს.

\paragraph{პოზიციური B+ ინდექსი რიცხვებში.}
ინდექსის აგებულება CRDT-ის თავში აღვწერეთ; აქ მის ეფექტზე ვისაუბრებთ.
ოპტიმიზაციამდე ყოველი კლავიში დოკუმენტის დაკავშირებული სიის თავიდან
ჩამოვლას მოითხოვდა, $O(n)$ ორივე მიმართულებით (პოზიციიდან ბლოკამდე და
ბლოკიდან პოზიციამდე). მცირე ფაილებზე ეს შეუმჩნეველია, მაგრამ ბლოკების
რაოდენობის ზრდასთან ერთად ღირებულება წრფივად იზრდება და ეს ღირებულება
\emph{ყოველ} კლავიშზე და \emph{ყოველ} დაშორებულ ოპერაციაზე გადაიხდება.
ჩვენი გაზომვებით (debug ბილდი, სადაც ხის განშტოების ფაქტორი განზრახ
მცირეა):

\begin{table}[htbp]
  \centering
  \begin{tabular}{rrrr}
    \toprule
    ბლოკების რაოდენობა & სიის ჩამოვლა ($O(n)$) & B+ ხე ($O(\log n)$) & აჩქარება \\
    \midrule
    $\sim$200    & $\sim$127~$\mu$s  & $\sim$14~$\mu$s & $9\times$ \\
    1\,000       & $\sim$620~$\mu$s  & $\sim$16~$\mu$s & $\sim 39\times$ \\
    10\,000      & $\sim$6\,200~$\mu$s & $\sim$18~$\mu$s & $\sim 340\times$ \\
    \bottomrule
  \end{tabular}
  \caption{პოზიციის ძებნის დრო ინდექსამდე და ინდექსის შემდეგ}
  \label{tab:btree}
\end{table}

საკვანძო დაკვირვება ისაა, რომ ხის სვეტი თითქმის მუდმივია, ლოგარითმი
პრაქტიკაში „ბრტყელია“. Release ბილდში (განშტოება 64/32, სრული
ოპტიმიზაციები) თითო ძებნა 1~$\mu$s-ზე ნაკლებია დოკუმენტის ზომის მიუხედავად,
ანუ რედაქტორის კადრის ბიუჯეტს ინდექსი ვერასდროს ემუქრება. ეს ზუსტად ის
მიზეზია, რის გამოც მოწიფული CRDT იმპლემენტაციები (Yjs, diamond-types
\cite{diamond}) მსგავს ინდექსებს აუცილებლობად თვლიან.

\paragraph{ნაგვის შემგროვებელი, რატომ არის საჭირო.}
როგორც აღვნიშნეთ, CRDT-ში წაშლა მხოლოდ მონიშვნაა და წაშლილი ბლოკები
tombstone-ებად რჩება. აქტიური რედაქტირებისას ეს სწრაფად გროვდება:
დიდხანს მიმდინარე სესიაში დოკუმენტის მეხსიერების უმეტესი ნაწილი შესაძლოა
დიდი ხნის წაშლილმა ტექსტმა დაიკავოს, რომელიც სნეფშოტებთან ერთად
ქსელშიც მოგზაურობს. tombstone-ის უბრალოდ წაშლა კი საშიშია: თუ რომელიმე
მონაწილემ ის ჯერ არ ნახა, ან თუ დაგვიანებული ოპერაცია მას origin-ად
ეყრდნობა, ნაადრევი წაშლა რეპლიკების უხმო დაშორებას (divergence)
გამოიწვევს. სწორედ ამიტომ ეთმობა Yjs-ში ნაგვის შეგროვებას ცალკე,
საკმაოდ ფაქიზი მექანიზმი.

\paragraph{ჩვენი მიდგომის თავისებურება.}
აქ ვისარგებლეთ ჩვენი არქიტექტურის ერთი თვისებით, რომელიც „სუფთა“
peer-to-peer სისტემებს არ აქვთ: ჩვენს სესიას ჰყავს გამორჩეული კვანძი -
ჰოსტი, რომელიც ყოველთვის ონლაინაა (სესია მის გარეშე ვერ არსებობს) და
რომელსაც გეითვეი წევრობის ზუსტ სურათს აწვდის. ეს ჰოსტს ბუნებრივ
კოორდინატორად აქცევს, ისე, რომ დოკუმენტის მონაცემები მაინც სრულად
დეცენტრალიზებული რჩება. სქემა ასე მუშაობს:

\begin{enumerate}[itemsep=0pt]
  \item ყოველი სტუმარი პერიოდულად (დებაუნსით, რომ ქსელი არ გადაიტვირთოს)
        აქვეყნებს თავის მდგომარეობის ვექტორს, „აი, რამდენი მაქვს
        ნანახი თითოეული ავტორისგან“;
  \item ჰოსტი აგროვებს ყველა წევრის ვექტორს (გეითვეის წევრობის ფრეიმები
        ეუბნება, ვინ არის ამჟამად ოთახში) და ითვლის მათ \emph{მინიმუმს} -
        „იატაკს“: ოპერაციების იმ ზღვარს, რომელიც ყველას გარანტირებულად
        აქვს ნანახი;
  \item ამ იატაკს ჰოსტი კვეთს საკუთარ DeleteSet-თან და იღებს
        „დადასტურებული“ წაშლების სიმრავლეს, წაშლებს, რომლებიც ყველამ
        ნახა;
  \item დადასტურებული დიაპაზონები \texttt{gc-commit} ფრეიმად ეგზავნება
        ყველა მონაწილეს (და ჰოსტისვე დოკუმენტს), რომლებიც მას
        ერთნაირად ასრულებენ.
\end{enumerate}

\texttt{gc-commit}-ის შესრულება განზრახ კონსერვატიულია: ბლოკები
სტრუქტურულად \emph{არ} იშლება, მხოლოდ მათი ტექსტური შიგთავსი იცლება და
სნეფშოტისთვის განკუთვნილი წაშლის სიმრავლეები იკვეცება. ბლოკის
იდენტიფიკატორი და ადგილი სიაში უცვლელი რჩება, ამიტომ დაგვიანებული
ოპერაცია, რომელიც ამ ბლოკს origin-ად იყენებს, ისევ ზუსტად ინტეგრირდება -
სწორედ ეს ხდის ჩვენს ნაგვის შეგროვებას კონვერგენციისთვის უსაფრთხოს.
სტრუქტურული წაშლა შემდეგი ნაბიჯი იქნებოდა, მაგრამ ის ინტეგრაციის
ალგორითმის ინვარიანტებს არღვევს (გაყოფის მიმართ ინვარიანტულობა დაიკარგება)
და მეხსიერების მოგება რისკად არ ღირს, მთავარი მოგება (ტექსტის
გათავისუფლება და სნეფშოტების შემცირება) ისედაც მიიღწევა.

რამდენადაც ჩვენთვის ცნობილია, ეს ჰოსტ-ავტორიტეტული,
მდგომარეობის-ვექტორის-იატაკზე დაფუძნებული სქემა განსხვავდება გავრცელებული
CRDT იმპლემენტაციების მიდგომებისგან, რომლებიც ან სრულად სიმეტრიულ
პროტოკოლებს იყენებენ, ან ნაგვის შეგროვებას საერთოდ დოკუმენტის „შენახვის“
მომენტამდე დებენ. ჩვენს შემთხვევაში ჰოსტის გამორჩეულობა ისედაც სისტემის
ნაწილი იყო, მისი გამოყენება ნაგვის შეგროვების გასამარტივებლად
არქიტექტურის ბუნებრივი გაგრძელება აღმოჩნდა.

\paragraph{პრობლემები, რომლებსაც გზაში გადავაწყდით.}
გაიდლაინის შესაბამისად, აქვე შევაჯამებთ განვითარებისას წამოჭრილ
მნიშვნელოვან პრობლემებსა და მათ გადაწყვეტებს:

\begin{itemize}[itemsep=2pt]
  \item \textbf{პოზიციების „ცურვა“ webview-სა და ბექენდს შორის} -
        ასინქრონული IPC-ის გამო ლოკალური ბრძანება შესაძლოა მოძველებულ
        პოზიციას ატარებდეს. გადაწყვეტა: \texttt{baseSeq} +
        ბექენდში დაშორებული ოპერაციების მცირე ჟურნალის მიმართ პოზიციის
        ტრანსფორმაცია.
  \item \textbf{Windows-ის CRLF დესინქრონიზაცია}, CRLF-იანი ფაილის
        გახსნისას პოზიციები პლატფორმებს შორის აღარ ემთხვეოდა და
        ხაზის შუაში რედაქტირება ცდებოდა. გადაწყვეტა: LF-ზე ნორმალიზაცია
        გახსნისას და CRLF-ის აღდგენა შენახვისას.
  \item \textbf{Monaco-ს readOnly ხაფანგი}, წაკითხვის რეჟიმში Monaco
        პროგრამულ რედაქტირებებსაც აგდებდა და სტუმარი დაშორებულ
        ცვლილებებს ვეღარ იღებდა. გადაწყვეტა: დროშის წამიერი მოხსნა
        დაშორებული ცვლილების ჩასმისას.
  \item \textbf{სნეფშოტ-სინქრონიზაციის რბოლა}, თავდაპირველად ჰოსტი
        სნეფშოტს მხოლოდ შეერთებისას აგზავნიდა და გვიან შემოსული სტუმარი
        შესაძლოა მოძველებულ მდგომარეობას იღებდა. გადაწყვეტა: მოთხოვნა-
        პასუხის სქემა, გეითვეი ყოველ ახალ წევრზე ჰოსტს \emph{ახალ}
        სნეფშოტს სთხოვს.
  \item \textbf{დიდი ჩასმების დანაწევრება}, გიგანტური ერთიანი ბლოკები
        (მაგ. მთელი ფაილის ჩასმა) ინდექსსა და ქსელს ტვირთავდა.
        გადაწყვეტა: ტექსტის 64-სიმბოლოიან ბლოკებად დაჭრა გახსნისას.
  \item \textbf{ინდექსისა და სიის დაშორება}, B+ ხის დანერგვისას
        ყველაზე სახიფათო რისკი ორი სტრუქტურის უხმო დაშორება იყო.
        გადაწყვეტა: debug-ორაკული, რომელიც ყოველი მუტაციის შემდეგ სრულ
        შედარებას აკეთებს, პლუს ხის შიდა ინვარიანტების ცალკე ვალიდატორი.
\end{itemize}
